\section{The Path-Finding Operator}
\label{sec:operator}

% Revised 27 Oct 2026 against paper/STYLE.md; numbers from eval-freeze-rc1 except the red laptop passages.

A path pattern in \sqlpgq asks, for a set of candidate (source, destination) pairs, whether a path of the allowed
length exists and how long it is, or which vertices and edges it uses. The operator takes exactly that. Its inputs are
a stream of pairs, one \duckdb vector of 2048 rows at a time, and the graph, either as a set of edges to build the CSR
from or as the name of a CSR index to read it from; its output is one length or one path per pair. Label tests,
filters on the source and destination vertices, the join back to the vertex tables and the test on the quantifier stay
relational and are left to \duckdb's planner. The operator is opt-in through the setting
\texttt{experimental\_path\_finding\_operator}; without it the rewrite emits the UDF calls of the CIDR design. Unless stated otherwise, the numbers in this section are taken from the evaluation of Section~\ref{sec:eval}
(LDBC SNB \texttt{Person\_knows\_Person}, one m7a.8xlarge, 32 threads). This section describes how the operator gets
its inputs, how it builds and lays out the graph, how it traverses it, and how it runs in parallel without any task blocking a thread while
it waits for another.

\subsection{From \texttt{MATCH} to an operator}
\label{sec:planner}

\duckpgq rewrites \texttt{MATCH} into SQL, and the path pattern becomes a scalar function call in that SQL. We call it
the \emph{marker}. It takes the source and destination columns of the candidate pairs, the edges as a struct column of
(source, destination) row ids, the vertex and edge counts, a constant key naming the graph, and the upper bound of the
quantifier. The marker is registered as a volatile scalar function so that no part of \duckdb folds or reorders it,
and it is never executed. An optimizer extension runs before \duckdb's own optimizer, finds the projection that
contains the marker, replaces the marker by a reference to the operator's output column, and places a logical
path-finding operator beneath the projection, with the pairs as its first child and, when the CSR is built in the
query, the edges as its second. The rule runs on the
binder's plan shape, before join reordering, because the two inputs are easy to identify there and hard to find after
the optimizer has pushed filters into them. \duckdb then sees an operator with one or two children and a cardinality, plans
the rest of the query around it, and executes it as a sink over its children and a source for the result.

Three details of the rewrite matter for performance. First, the filters on the source and destination vertices are
applied to each vertex set before the two are combined, so a query with one source and a selective filter on the
destination yields few candidate pairs, and the unfiltered cross product of the vertex tables is never formed.
Second, the upper bound of the quantifier travels with the marker, so the traversal stops at that depth instead of
exploring further paths that a filter would discard afterwards. Third, when a CSR index exists for the edge table at bind time,
the rewrite drops the edge input and passes the index name in the key. The operator resolves the key to a snapshot
(Section~\ref{sec:index}), and the pairs carry vertex keys instead of row ids, because the index identifies vertices
by key.

\subsection{Building the partitioned CSR on the fly}
\label{sec:csr}

Without an index the operator builds the CSR on the fly, inside every query, as the CIDR design did. The build is cheap, but
when the traversal is cheaper still, the build takes most of the query. At SF300 a query for one pair takes 274\,ms with
the build and 10\,ms over the index (Table~\ref{tab:random-pairs}), so the build is 96\% of it. For 1024 random
pairs, where the traversal does more work, the build is two fifths of the query (1.3\,s against 0.77\,s). Figure~\ref{fig:partitioned-csr} shows the build. The
edges arrive as (source, destination) row-id pairs from a join of the edge table with its two vertex tables. The sink
packs each edge into 64 bits and scatters it into a bucket selected by the high bits of its destination id, similar to
\duckdb's hash join partitioning. A range spans at most 65,535 destinations, so that a destination fits in 16 bits as an offset into its range, and
there are at least eight ranges, so that a small graph still gives eight tasks; the width is a setting and
Section~\ref{sec:eval-ablation} measures it. Tasks then take the ranges largest first, sort each range's edges by
source and emit the partition of Figure~\ref{fig:partitioned-csr}b. The bucket index lets a lookup of one source read
one bucket entry and scan about two rows, where a binary search would take a logarithmic number of probes into the
rows of the partition. No per-vertex array exists at any point,
so the build memory is proportional to the edges and not to the vertices.

\begin{figure}[t]
  \centering
  % Figure: the partitioned CSR, (a) its build on the fly, (b) one partition. Column width (8.2 cm).
% Colours as in Figures 1 and 2: operator steps blue, data grey, the part the index adds orange.
\begin{tikzpicture}[
  font=\tiny,
  >={Stealth[length=3pt,width=2.6pt]},
  box/.style={draw, rounded corners=1pt, align=center, inner sep=1.5pt, minimum height=4mm, anchor=north west},
  data/.style={box, draw=black!55, fill=black!4},
  opstep/.style={box, draw=opblue, fill=opblue!10},
  fld/.style={draw=black!60, minimum height=4mm, inner sep=0pt, anchor=north west, align=center},
  cell/.style={draw=black!55, fill=white, minimum width=4.6mm, minimum height=3.6mm, inner sep=0pt, anchor=north west},
  rid/.style={cell, draw=ixorange, fill=ixorange!10},
  lbl/.style={font=\tiny, text=black!60, inner sep=0.5pt},
  head/.style={font=\scriptsize\bfseries, text=black!80, inner sep=0pt, anchor=north west},
]
% ------------------------------------------------------------------ (a) build on the fly
\node[head] at (0,0) {(a) building it on the fly};
\node[data, text width=2.0cm] (stream) at (0,-0.4) {edge stream\\ (src, dst) row ids};
\node[lbl, anchor=south west] at (2.75,-0.38) {one edge packed into 64 bits};
\node[fld, minimum width=1.3cm, fill=opblue!15] (f1) at (2.75,-0.42) {dst range};
\node[fld, minimum width=1.6cm, fill=white] (f2) at (f1.north east) {source id};
\node[fld, minimum width=1.25cm, fill=white] (f3) at (f2.north east) {local dst};
\draw[->, semithick, black!70] (stream.east |- f1.west) -- (f1.west);
% buckets
\foreach \i/\t in {0/range 0,1/range 1,2/range 2,3/$\cdots$,4/range $p$} {
  \node[data, minimum width=0.95cm] (b\i) at (2.2+\i*1.15,-1.25) {\t};
}
\draw[->, semithick, opblue] (f1.south) -- node[lbl, right, pos=0.45] {high bits select the bucket (sink)} (f1.south |- b1.north);
% tasks take the ranges largest first
\node[opstep, text width=7.3cm] (task) at (0.3,-1.95) {per range, largest first: radix sort by source $\rightarrow$ emit sparse rows $\rightarrow$ attach bucket index};
\foreach \i in {0,1,2,4} {\draw[->, opblue] (b\i.south) -- (b\i.south |- task.north);}

% ------------------------------------------------------------------ (b) one partition
\node[head] at (0,-2.75) {(b) one partition, destination range $p$};
\def\xo{1.75}      % x of the first edge cell
\def\cw{0.5}       % cell pitch
% bucket index above the sources: source >> 3 selects the bucket
\node[lbl, anchor=east] at (\xo-0.1,-3.3) {bucket index};
\foreach \k/\x in {0/0.25,1/1.5,2/3.05} {\node[cell, minimum width=4.2mm] (k\k) at (\xo+\x,-3.12) {$b_\k$};}
\node[lbl, anchor=west] at (\xo+3.65,-3.3) {$b_k$: first row with source $\gg 3 = k$};
% sparse rows: sources spanning their edge groups
\node[lbl, anchor=east] at (\xo-0.1,-3.95) {source};
\foreach \i/\v/\o/\n in {0/3/0/2,1/7/2/1,2/12/3/3,3/18/6/1} {
  \node[cell, minimum width=\n*\cw cm-0.4mm, fill=black!6] (s\i) at (\xo+\o*\cw,-3.77) {\v};
}
\draw[->, black!60] (k0.south) -- (s0.north);
\draw[->, black!60] (k1.south) -- (s2.north);
\draw[->, black!60] (k2.south) -- (s3.north);
\node[lbl, anchor=east] at (\xo-0.1,-4.38) {offset};
\foreach \o in {0,2,3,6,7} {\node[lbl] at (\xo+\o*\cw,-4.38) {\o};}
% edges: 16-bit local destinations and row ids
\node[lbl, anchor=east] at (\xo-0.1,-4.78) {local dst (16 bit)};
\foreach \j/\d in {0/5,1/9,2/1,3/0,4/4,5/8,6/2} {\node[cell] (e\j) at (\xo+\j*\cw,-4.6) {\d};}
\node[lbl, anchor=east] at (\xo-0.1,-5.2) {edge row id};
\foreach \j in {0,...,6} {\node[rid] at (\xo+\j*\cw,-5.02) {$r_{\j}$};}
\node[lbl, anchor=west, align=left] at (\xo+3.65,-4.62) {only sources with an edge into\\ range $p$ have a row};
\node[lbl, anchor=west, align=left] at (\xo+3.65,-5.13) {\textcolor{ixorange}{kept by the index, and by the}\\ \textcolor{ixorange}{build when paths are asked for}};
\end{tikzpicture}

  \caption{The partitioned CSR. (a) The on-the-fly build: the sink packs each edge into 64 bits and the high bits,
  the destination range, select its bucket; tasks take the buckets largest first, sort each by source and emit its
  partition.
  (b) One partition. Rows exist only for sources with an edge into the range; a source's edges are the slice between
  its offset and the next, stored as 16-bit offsets into the range. The bucket index maps the high bits of a source id
  to its first row. The row ids (orange) are what the CSR index adds, and what the build adds when a query asks for
  paths.}
  \label{fig:partitioned-csr}
\end{figure}

We call this layout the \emph{partitioned CSR}. Partitioning the adjacency by destination range so that a traversal
step writes only a cache-sized slice of the frontier is an established idea, known as CSR
segmenting~\cite{zhang2017cagra} and propagation blocking~\cite{beamer2017blocking}, and GraphChi's shards partition
edges by destination interval for the same reason~\cite{kyrola2012graphchi}. What is specific here is that the layout
is produced in one pass from a join stream inside a relational operator, with radix partitioning and no per-vertex
array, and that the same layout is the on-disk format of a transactional index. The operator's build and the index of
Section~\ref{sec:index} share it for a reason. A partition holds every edge into one destination range, so a traversal
step that processes partition $p$ writes only to vertices of range $p$, and the tasks of one level never write the
same vertex. Destinations inside a
partition are 16-bit offsets, half the size of 32-bit ids, and an index epoch stores exactly this layout plus the
row ids. The difference between the two routes for a single pair puts the build at 22, 128 and 264\,ms for the
1.84, 18.7 and 55.7\,M edges of SF10, SF100 and SF300. \todo{The CSR's size in bytes per edge is not measured in the
frozen run; measure it or drop the size claim.}

\subsection{Traversals}
\label{sec:traversals}

The operator has two traversals and a rule that picks one; the multi-source traversal assigns its lanes in one of
two ways. Both traversals run breadth-first over the partitions and both honour the depth bound. They differ in what a
step costs.

\paragraph{Multi-source traversal.} For many pairs we keep \msbfs~\cite{then2014msbfs}. 512 searches share one sweep
over the adjacency, each search owns a bit in three bit-set arrays per vertex (seen, this frontier, next frontier),
and relaxing an edge is a 512-bit \texttt{or}. A lane is either one pair or one source. With a lane per pair, each
input vector of up to 2048 pairs becomes one batch, in the order the vectors arrive, and a batch runs as waves of 512
lanes. With a lane per source, the formulation of Then et al., a wave packs up to 512 distinct sources, each with the
list of destinations its pairs ask for. After every level a lane records the level of the destinations it has newly
reached and drops them from its list, and a lane whose list is empty is masked out of the next frontier, so it stops
spreading. LDBC BI19 at SF300 asks for 2.4\,M pairs from 1.4\,K sources; that is three waves with a lane per source
and 4.6\,K waves with a lane per pair. A wave holds its three bit sets per vertex, 192 bytes at 512 lanes, which is
86\,MB at SF100. The batches of a lane-per-pair query run in worker groups, at most one per thread. When the working sets of
all groups fit a budget of 512\,MiB, all groups start at once and the bit sets stay in cache; otherwise the scheduler
starts 16 groups, doubles their number while the throughput rises by more than a tenth and steps back when it falls by
more than a tenth. A group whose working set is at most a sixteenth of the budget has its own threads; larger groups
pool their tasks, so a group that waits at the end of its level leaves no thread idle. The waves of a lane-per-source
query run side by side, each with its own working set. \todo{They bypass the worker groups and their memory budget:
a query with many waves allocates all their working sets at once. Route them through the scheduler, or say so.}

A level has two phases, an explore phase that relaxes the frontier into the next bit sets and a check phase that folds
the next bit sets into the seen bits. A summary bitmap marks the vertices with a live lane, and a level whose live
vertices are fewer than a sixteenth of all vertices visits only those, through the bucket index. A dense level sweeps
all sparse rows of all partitions and costs the same whether one lane or five hundred are alive. That is a good trade
when hundreds of lanes are alive and a bad one when one is. One pair over the index takes 13\,ms on the lanes against
4\,ms on the single-source traversal at SF10, and 45 against 9\,ms at SF100. \todo{Single runs before rc5; the rerun's
Section 3 microbenchmarks replace them.}

\paragraph{Single-source traversal.} For few sources we run one search per source with a frontier list, in the shape
Ligra~\cite{shun2013ligra} made standard. The frontier is a vector of vertex ids. While it is smaller than a
sixty-fourth of the vertex count, each frontier vertex is looked up in every partition through the bucket index and
its edges are relaxed with a compare-and-swap on a byte per vertex, so a step costs the frontier plus its edges. When
the frontier grows past that bound the step switches to a sweep over the partitions with a per-vertex level stamp
marking the frontier. The search records the destinations its group of pairs asks for and stops as soon as the last
one is reached, or at the depth bound. When the query asks for paths, the traversal also records, per vertex, the
vertex and the edge through which it was first reached, and the result scatter walks that chain backwards
(Section~\ref{sec:paths}). The multi-source traversal cannot afford that; a parent per lane per vertex is 512 times
the memory.

\paragraph{Choosing.} When the pairs are complete, the operator estimates the distinct sources and destinations with
HyperLogLog sketches (64 registers, about 13\% error, enough for a threshold). With at most 16 distinct sources the
single-source traversal runs, one group per source and the threads divided over the groups. With more sources and at
least four pairs per source the multi-source traversal runs with a lane per source, and otherwise with a lane per
pair. Paths and reach always use the single-source traversal, because a path needs a parent per vertex and reach
returns a distance per vertex, and the lanes keep neither. The same sketches decide the orientation. When
there are at least four times as many distinct sources as destinations and reverse partitions are available, which is
the case for an undirected pattern and for an index that stores its reverse side (Section~\ref{sec:index}), the
search starts from the destinations; the on-the-fly build produces forward partitions only. Both rules are
thresholds today, and Section~\ref{sec:strategy} discusses what a cost model needs in order to replace them. The
reverse orientation pays off when many sources share one destination, the shape of LDBC BI20
(Section~\ref{sec:eval-ablation}). \todo{Confirm after the full sweep whether the reverse side stays; then describe it
in Sections 3.2 and 4.} Twenty sources with 200 destinations each, over the index, take 13\,ms on the single-source
traversal, 20\,ms with a lane per source and 86\,ms with a lane per pair at SF10, and 91, 79 and 509\,ms at SF100.
On BI19 at SF300 the lanes per source take 1.6\,s against 30.6\,s for one search per source, and 4.7 against
43.8\,GiB of memory at the peak, because each single-source search keeps a distance, a seen byte and a level stamp
per vertex and 1.4\,K of them run at once. \todo{Single runs before rc5; confirm from the rerun.}

\paragraph{Unbound destinations.} A pattern whose destination carries no condition, every person reachable from $a$,
would produce one candidate pair per vertex. The rewrite recognises this case and asks the operator for the
\emph{reach} of each source instead, one single-source traversal per source that returns every reached vertex with its
distance; the destination table is joined back afterwards only when the query projects its columns. The sources
$\times$ all queries of Section~\ref{sec:eval} run this way.

\paragraph{Pairs with an unknown source or destination.} A pair whose source or destination is \texttt{NULL}, or whose
key is absent from the index, has no path. The operator answers \texttt{NULL} for it without searching and without
changing the traversal chosen for the other pairs of the query. A vertex without edges is searched like any other and
its search ends after one level.

\subsection{Parallelism as a chain of events}
\label{sec:phases}

A BFS level is a parallel sweep followed by a serial decision about which vertices were reached, whether any lane is
still alive and whether the depth bound is hit. The operator runs every traversal as a \emph{chain of per-phase
events}. A phase is a set of tasks that run independently and finish; when the last one finishes, the event's
completion callback runs the serial step once and inserts the event of the next phase. The on-the-fly build has the
same shape with fewer phases. The partitioning happens in the sink as the edges stream in, one event's tasks sort and
emit the ranges, and its completion callback publishes them. Each phase still waits for the one before it, but no task ever blocks its thread waiting for another task. This is the shape
\duckdb's own operators have, and it is what lets the operator share \duckdb's fixed thread pool with other queries.
The scheduler may interleave the tasks of concurrent queries freely, and an error or an interrupt stops the operator
at the next task boundary. The same abstraction covers the multi-source traversal
(two phases per level and batch) and the single-source traversal (one event per level and source group; a group with a
single worker runs all its levels in one task). The cost is one or two events per level, which is microseconds. The
concurrent-reader experiment of Section~\ref{sec:eval} runs 16 readers against one writer at 32 workers per query
with no coordination beyond \duckdb's scheduler.

\subsection{Paths}
\label{sec:paths}

\texttt{element\_id(p)}, \texttt{vertices(p)} and \texttt{edges(p)} are answered by the same operator in a second
mode that returns, per pair, the list $[v_0, e_0, v_1, \dots, v_k]$ of vertices and edge row ids. The vertices are
row ids when the CSR is built in the query and keys when it is read from an index. The
single-source traversal records parents as described above; the scatter walks from the destination to the source,
maps vertex ids to row ids or keys and edge positions to row ids, and reverses the list. \texttt{path\_length(p)} is then
$\lfloor\text{len}/2\rfloor$, and \texttt{vertices} and \texttt{edges} are the even and odd positions, which is also
how the SQL fallback of \duckpgq defines them, so the build route and the fallback return the same lists. Edge row ids are available because
the build keeps one per edge position when paths are requested and the index keeps them always
(Section~\ref{sec:index}), so a path can be joined back to the edge table for its properties. Paths cost more than
lengths when the pairs have many sources, because paths need the single-source traversal and lengths use the lanes.
For 2048 random pairs over the index, paths take 3.6 times as long as lengths at SF10 and 4.5 times at SF100.
\todo{Single runs before rc5; confirm from the rerun.}
